% chapters/06/sub-cpmk.tex
\begin{subcpmkbox}
Bab ini menuntaskan \SUBCPMK{1.3} secara utuh. Ketiga indikatornya diperiksa,
dan buktinya berupa prototipe permainan yang benda-bendanya jatuh, melompat,
dan memantul.

\vspace{0.5em}
\textbf{\SUBCPMK{1.3} Menerapkan logika fisika sederhana (gravitasi, lompatan,
pantulan) dalam simulasi 2D} --- ranah C3, \CPMK{1}, \CPL{03}.

\vspace{0.5em}
{\setlength{\tabcolsep}{4pt}\footnotesize%
\begin{tabularx}{\textwidth}{@{}C{1.9cm} L{4.0cm} L{2.5cm} Y@{}}
\toprule
\headrow \textbf{Indikator} & \textbf{Rumusan} & \textbf{Dibahas pada} & \textbf{Bukti} \\
\midrule
\ind{a} & Menerapkan integrasi Euler untuk memperbarui kecepatan dan posisi &
Bagian~\ref{sec:percepatan} &
Kode \code{ch06_fisika.py} pada repositori, khususnya urutan pembaruan
kecepatan lalu posisi. \\
\addlinespace
\ind{b} & Mengimplementasikan lompatan yang hanya aktif ketika objek berada di
tanah &
Bagian~\ref{sec:lompatan} &
Kode lompatan beserta uji bahwa menahan tombol tidak menghasilkan lompatan
berulang. \\
\addlinespace
\ind{c} & Menghasilkan pantulan dengan koefisien restitusi sederhana &
Bagian~\ref{sec:pantulan} &
Kode pantulan pada lantai dan dinding, beserta catatan pengukuran tinggi
pantulan pada lembar kerja praktikum. \\
\bottomrule
\end{tabularx}}
\end{subcpmkbox}
